\section{Motivation}
\label{sec:benchmarks_motivation}

Chapter~\ref{ch:foundations} defined open-world environments as settings in which an agent's internal models are incomplete with respect to the environment, and formalized novelty as a mismatch between those models and observed outcomes. In such settings, the central question is not only how agents should adapt, but how their ability to detect, localize, and recover from mismatch can be evaluated in a principled and systematic manner.

Evaluating agents under mismatch presents challenges that do not arise in closed-world settings. In standard benchmarks, the environment is fixed and fully specified, and performance can be measured with respect to a known objective. In contrast, open-world environments are characterized by the possibility of unanticipated changes: objects may appear or disappear, action effects may change, and task structure may evolve. These changes induce mismatch between the agent's expectations and the environment, but the space of possible mismatches is unbounded. As a result, evaluation cannot rely on a fixed set of scenarios, but must instead provide mechanisms for generating controlled and interpretable variations in the environment.

The crafting scenario introduced in Chapter~\ref{ch:foundations} illustrates this challenge. When the agent encounters the axe novelty---where trees become unbreakable without a new tool---its existing plan fails due to a mismatch between the expected and actual effects of the \texttt{break} operator. Evaluating the agent's behavior in this setting requires more than measuring task success: it requires quantifying how quickly the agent detects the mismatch, how severely its performance degrades, and how efficiently it recovers. More generally, evaluation must capture the dynamics of adaptation, not just its final outcome.

Moreover, evaluating the central claim of this dissertation---that structured abstractions over actions and interactions determine whether agents can adapt efficiently under novelty (\cref{sec:abstractions})---requires benchmarks that can distinguish between agents with different representational commitments. An environment that only measures final task success cannot reveal whether a hybrid agent's advantage over a monolithic learner stems from its structured decomposition of behavior or from some other factor. The benchmarks introduced in this chapter are designed to expose these differences through metrics that capture adaptation speed, performance degradation, and the ability to exploit beneficial changes.

To evaluate open-world agents along these dimensions, we require environments that can systematically introduce mismatch, expose its effects at both the planning and execution levels, and support agents with different representational and decision-making paradigms. This chapter presents two such environments. NovelGridworlds introduces a structured taxonomy of novelties and evaluation metrics that capture performance degradation and detection. NovelGym extends this framework with a formal environment representation, composable transformations that induce mismatch, and metrics that explicitly characterize adaptation dynamics. Together, these environments provide the experimental foundation for studying adaptation and generalization under novelty in the remainder of this dissertation.